\begin{frame}[fragile]{Ponto flutuante: valores aproximados}
\begin{lstlisting}
float temperatura_c = 27.5f;
double tensao_v = 3.3;
\end{lstlisting}
\begin{itemize}
\item No código C, a parte decimal usa \textbf{ponto}.
\item O sufixo \texttt{f} especifica um literal \texttt{float}.
\item Sem o sufixo, \texttt{3.3} tem tipo \texttt{double}.
\item Muitos valores decimais têm representação binária aproximada; não suponha precisão infinita.
\end{itemize}
\begin{block}{Unidade no nome}
\texttt{temperatura\_c} e \texttt{tensao\_v} ajudam a evitar confusões entre grandezas e unidades.
\end{block}
\end{frame}
